% !PRINTED 296
\weberpage{296}
%<<CONT
\gesperrt{2. Wenn eine irreducible Gleichung, in deren Grad mehr als eine
Primzahl aufgeht, durch successive Adjunction von Radicalen in Factoren
zerfällt, so wird eine Zerfällung in $s$ Factoren $r^{\mathrm{ten}}$ Grades
herbeigeführt durch Adjunction der Wurzeln einer metacyklischen Gleichung
$r^{\mathrm{ten}}$ Grades.}

Dieser Satz enthält als speciellen Fall den von \gesperrt{Abel} herrührenden
Satz, dass eine irreducible Gleichung, deren Grad $m$ nicht die Potenz einer
Primzahl ist, nur dann durch Radicale lösbar sein kann, wenn sie durch
Adjunction der Wurzeln einer lösbaren Gleichung niedrigeren Grades, deren Grad
ein Theiler von $m$ ist, in Factoren zerfällt\textsuperscript{1)}. Damit ist die
Frage nach der Auflösung einer Gleichung durch Radicale oder auch nur der
Reduction einer Gleichung durch Radicale ausserordentlich vereinfacht. Man kann
in der That die fernere Untersuchung auf Gleichungen beschränken, deren Grad
eine Primzahl oder eine Primzahlpotenz ist, weil darauf alle anderen Fälle durch
wiederholte Anwendung des Satzes 2. zurückgeführt sind.

So gestattet z.\,B. die Frage nach allen durch Radicale lösbaren irreduciblen
Gleichungen $6^{\mathrm{ten}}$ Grades eine geradezu triviale Antwort:

\gesperrt{Um alle metacyklischen Gleichungen $6^{\mathrm{ten}}$ Grades in irgend einem
Körper $\Omega$ zu erhalten, adjungire man dem Körper $\Omega$ eine
Quadratwurzel und bilde in dem erweiterten Körper alle cubischen Gleichungen,
oder man adjungire die Wurzel einer cubischen Gleichung und bilde in dem
erweiterten Körper alle quadratischen Gleichungen.}

Als Beispiel einer solchen Gleichung führen wir die von \gesperrt{Hesse}
behandelte an, von der die Kreisschnitte einer nicht auf die Hauptaxen
bezogenen Fläche $2^{\mathrm{ten}}$ Grades abhängen. Dies Problem wird durch
Quadratwurzeln gelöst, wenn vorher durch die Lösung einer cubischen Gleichung
die Hauptaxen bestimmt sind\textsuperscript{2)}.

\noindent\rule{4cm}{0.4pt}\par
\textsuperscript{1)} \gesperrt{Abel} giebt den Satz ohne Beweis in der
Abhandlung \glqq Sur la rés. algébr. des équations\grqq{}. Oeuvres complètes
1881, Bd. II, S. 217. Vgl. auch die Abhandlung von \gesperrt{Galois} in
\gesperrt{Liouville's} Journal, Bd. 11.\par
\textsuperscript{2)} \gesperrt{Hesse}, \glqq Ueber die Auflösung derjenigen
Gleichungen $6^{\mathrm{ten}}$ Grades etc.\grqq{} Crelle's Journal, Bd. 41
(1851).
